\section{Introduction}\label{sec:intro}

A smart account that authenticates with a passkey is lost when every device
holding the passkey is lost. Account-abstraction wallets answer this with
recovery: guardians who approve a key change, a delay during which the owner
can object, and secondary credentials~\cite{buterin-social-recovery,argent-spec,erc4337}.
An email inbox is the secondary credential most people already have and keep
for years, and large services have long used it to recover
accounts~\cite{bonneau-recovery}. Used naively, it puts the address in the
operator's database and makes whoever runs the email check a party that can
recover any account: a server that mails a one-time code, or a recovery
server that holds signing keys for its users~\cite{sep30}, is trusted with
every account it serves. Checking the DKIM signature~\cite{rfc6376} on chain
instead removes the server but publishes the signed \texttt{From} header.

This paper shows that control of an inbox can be one recovery factor while
the chain holds only a salted hash of the address,
$\Commit$, where $\code$ is an account code that only the user holds, and
the operator's servers see the address only inside an attested enclave. To
recover, the user sends one email from the inbox to a fixed address. Its
Subject carries a blinded copy of the digest of the pending recovery
statement, so the email's content names no account or attempt. The user's
browser proves in zero knowledge that the provider DKIM-signed the email
under a registered key, that its sender opens the commitment, and that its
Subject opens to the statement's digest. A per-account email-guardian
contract verifies the proof on Soroban and counts it as one guardian
approval in Perch's recovery~\cite{perch-recovery-spec}, used unchanged.
Perch's quorum and delay still apply: email approvals never complete a
recovery before the delay has passed, and they authorize one alone only
under a lost-key profile with at least seven days of delay and a second
alert channel, where the owner's cancellation is the defence.

Five problems shape the design. \emph{Binding}: an email must approve exactly
one recovery statement (one account, new key, attempt and deadline), and a
copy must be worth no more than the original. \emph{Parsing soundness}: a
circuit must read a header the way the provider signed it; public audits of
open-source email circuits have repeatedly found sender spoofing and
header-confusion
flaws~\cite{audit-zksecurity-2024,audit-yacademy,audit-veridise-2024,audit-consensys-2024,audit-matterlabs-2024,audit-zellic-2024}.
\emph{Key trust}: whoever says which RSA key signs for \texttt{gmail.com}
must not be able to plant a key and recover accounts with it. \emph{Privacy
from the operator}: the operator's servers and staff must not read the
address at enrollment, at recovery, or when sending alerts. \emph{The
human}: a user can be talked into sending the email and reading out the
code, and under email plus delay only an alert that reaches the owner in
time stops a wrongful recovery.

\paragraph{Contributions.}
\begin{itemize}
  \item A design for email as one guardian inside an existing recovery
    protocol, unchanged (\S\ref{sec:design}). What is new against zkEmail's
    recovery stack (\S\ref{sec:related}): the Subject binds the protocol's
    own statement digest, blinded, rather than a readable command; the proof
    is made in the user's browser, so no relayer receives the email or holds
    the account code; the key registry is immutable, fed by DNS attestations,
    with public timelocks, and the per-account pin is salted, expiring and
    fixed with the guardian; and alerts come from the protocol's own events
    through a sealed vault, with a liveness probe and a wallet that reads
    open attempts itself\claim{guardian-on-perch}.
  \item Four properties stated over a Lean model of the guardian, the
    circuit's relation and Perch's controller (\S\ref{sec:properties}).
    \emph{Proven} over the model: soundness and replay resistance, each the
    property or a hash collision on inputs the theorem computes, with
    knowledge extraction per accepted proof; one-factor safety over every
    trace; privacy as a simulation theorem in an ideal model of the commitment
    and the
    proofs\claims{email-factor-soundness,replay-resistance,one-factor-safety,address-privacy}.
    \emph{Assumed}: the cryptography behind the ideal model and the named
    breaks, and that providers authenticate senders. \emph{Tested}: the
    model's agreement with the contracts, and the circuit's parsing.
  \item A Noir circuit that refuses each publicly documented soundness issue
    class of open-source email circuits at a named assertion, and whose
    parsing agrees with reference implementations on real mail and fuzzed
    headers\claims{circuit-issue-classes,parser-real-mail,parser-fuzzing}.
  \item An evaluation on the test host, a local network running mainnet's
    settings, and testnet (\S\ref{sec:eval}): a 73-row attack suite that
    includes phishing, deleted alerts, a lapsed push, a watcher outage and
    repeated attempts; each environment's fees attributed to its settings;
    archival and its keeper; proving time on a desktop; phones not measured.
  \item A claim ledger giving every claim a status and its evidence, checked
    by the build (Appendix~\ref{app:claims}).
\end{itemize}

\paragraph{Claim tags.} A tag ends each statement about the system:
\claimtag{proven} for a Lean theorem about the model, \claimtag{tested} for
named tests, \claimtag{measured} for a recorded measurement,
\claimtag{assumed} for a stated assumption, \claimtag{unproven} for a
statement without evidence. Appendix~\ref{app:claims} lists each claim with
its evidence; the build fails if a tag names no ledger entry, if evidence is
missing, or if a measured figure does not occur in its measurement files.
